\documentclass[10pt,a4paper]{amsart}
\usepackage[margin=19mm]{geometry}
\usepackage{amsmath,amssymb,mathtools}
\usepackage[scr=rsfs,cal=euler]{mathalfa}
\usepackage{xfrac}
\usepackage[unicode,hidelinks,bookmarks=false]{hyperref}
\newcommand{\ie}{i.e.\ }
\newcommand{\cf}{cf.\ }
\newcommand{\resp}{respectively\ }
\newcommand{\Subsection}{\subsection}
\providecommand{\nobreakdash}{\nobreak\hbox{-}}
\setlength{\emergencystretch}{2em}
\multlinegap0pt
\usepackage{amssymb}
\usepackage{mathtools}
\usepackage{dsfont}
\usepackage{xcolor}
\newtheorem{theorem}{Theorem}
\def\N{{\mathbb N}}
\def\R{{\mathbb R}}
\newcommand{\Rd}{{\R^d}}
\def\T{{\mathbb T}}
\title{On the existence problem of regular \\ Gabor frames}
\author{Jaume de Dios Pont, Lukas Liehr, Mitchell A. Taylor}
\date{Source: arXiv:2606.26052v1 (CC BY 4.0)}
\begin{document}
\maketitle

\noindent\emph{Source and licence.} On the existence problem of regular \\ Gabor frames. Version arXiv:2606.26052v1 (CC BY 4.0), distributed by arXiv under the Creative Commons Attribution 4.0 International licence, http://creativecommons.org/licenses/by/4.0/, as recorded in the arXiv licence metadata for this exact version and shown on the abstract page https://arxiv.org/abs/2606.26052v1. Full licence notice: \texttt{licenses/arxiv-2606.26052-CC-BY-4.0.txt}.\par\smallskip

\noindent\textit{Editorial scope.} Counted unit: the article's theorem thm:zak\_zero (source0.tex lines 227--238). The article's complete original proof of it is delivered with it (source0.tex lines 507--809, one \texttt{proof} environment of 303 lines, which the article places away from the statement). No mathematics is written, repaired or re-derived: the statement and the proof are the article's own lines, in the article's own notation, and no retained line is substituted or re-flowed.  No further source-local context is needed: every cross-reference of every retained line resolves inside the two delivered spans.  Nothing was omitted from any retained span, so the package carries no omission record.  No retained line contains a citation, so the wrapper carries no bibliography and no bibliography entry.  No retained line contains a figure environment, an image-inclusion command or drawing code, so no image is delivered and the package declares no asset dependency.  Editorial formatting only, in addition: an AMS \texttt{amsart} preamble that restates the article's own declarations and macros used by the retained lines, namely the environments \texttt{theorem} on separate counters, together with the standard packages those lines need.  The retained environments print in this document as Theorem 1 in order of appearance, whereas the article prints the same items with the numbering of its own full document.  The complete native source of the article as distributed in the arXiv e-print for this exact version is bundled unchanged as \texttt{sources/203664/source0.tex}.
\begin{theorem}\label{thm:zak_zero}
Let $d,N\in \N$, and suppose that
$$
s_1,\dots,s_N:\mathbb R^d \times \mathbb R^d \to \mathbb C
$$
are continuous Zak-quasiperiodic functions. If $N \leq d$, then $s_1,\dots,s_N$ have a common zero, i.e., there exists $(u_0,v_0)\in\mathbb R^d\times\mathbb R^d$ such that
$$
s_1(u_0,v_0)=s_2(u_0,v_0)=\cdots=s_N(u_0,v_0)=0.
$$
If $N>d$, then the conclusion is false: there exist smooth Zak-quasiperiodic functions
$s_1,\dots,s_N :\mathbb R^d \times \mathbb R^d  \to \mathbb C$ (not all equal to zero) such that $s_1,\dots,s_N$ do not have a common zero.
\end{theorem}

\begin{proof}[Proof of the first part of Theorem \ref{thm:zak_zero}]    

In the following, we denote an element of $\Rd \times \Rd$ by $(u,v)$. Assume by contradiction that
\begin{equation}\label{eq:star}
    (s_1(u,v),\dots,s_N(u,v))\neq 0, \quad (u,v)\in\mathbb R^d\times\mathbb R^d.
\end{equation}

\smallskip
\noindent\textbf{Step 1: the quasiperiodic line bundle.}
Let
$
\mathbb T^{2d}
:=
(\mathbb R^d/\mathbb Z^d)_u\times(\mathbb R^d/\mathbb Z^d)_v
$
and define a complex line bundle
$
\mathcal L \to \mathbb T^{2d}
$
as follows. Start with the trivial bundle
$
\mathbb R^{2d}\times \mathbb C\to \mathbb R^{2d}
$
and let $\mathbb Z^d\times\mathbb Z^d$ act on $\mathbb R^{2d}\times \mathbb C$ by
$$
(n,m)\cdot (u,v,z)
:=
\bigl(u+n,\ v+m,\ e^{2\pi i\,n\cdot v}z\bigr),
$$
where $n,m\in\mathbb Z^d$.
This is a well-defined group action since for every pair $(n,m),(n',m')\in
\mathbb Z^d\times\mathbb Z^d$, one has
$$
\begin{aligned}
(n,m)\cdot\bigl((n',m')\cdot(u,v,z)\bigr)
&=
(n,m)\cdot
\bigl(u+n',v+m',e^{2\pi i\,n'\cdot v}z\bigr)  \\
&=
\bigl(u+n+n',v+m+m',
e^{2\pi i\,n\cdot(v+m')}e^{2\pi i\,n'\cdot v}z\bigr)  \\
&=
\bigl(u+n+n',v+m+m',
e^{2\pi i\,(n+n')\cdot v}z\bigr) \\
&=
(n+n',m+m')\cdot(u,v,z).
\end{aligned}
$$
Thus, we may form the quotient bundle
$$
\mathcal L
:=
(\mathbb R^{2d}\times\mathbb C)/(\mathbb Z^d\times\mathbb Z^d)
\to
\mathbb T^{2d}.
$$
A continuous section of $\mathcal L$ is equivalently a continuous function
$
\sigma:\mathbb R^{2d}\to\mathbb C
$
satisfying
\begin{equation}\label{eq:333}
    \sigma(u+n,v+m)
=
e^{2\pi i\,n\cdot v}\sigma(u,v)
\end{equation}
for all $n,m\in\mathbb Z^d$.
Indeed, such a function defines a section by
$$
[u,v]\longmapsto [(u,v,\sigma(u,v))],
$$
and condition \eqref{eq:333} is exactly the condition that this definition is independent of
the chosen representative $(u,v)$ of the point $[u,v]\in\mathbb T^{2d}$.
By assumption, each $s_\ell$ satisfies \eqref{eq:333}. Hence, each $s_\ell$ defines a continuous
section of $\mathcal L$ and therefore
$$
S:=(s_1,\dots,s_N)
$$
defines a continuous section of the rank-$N$ complex vector bundle
$$
E_N:=\mathcal L^{\oplus N} \to \mathbb T^{2d}.
$$
Assumption \eqref{eq:star} says exactly that $S$ is a nowhere-zero section of
$E_N$.

\smallskip
\noindent\textbf{Step 2: computation of the first Chern class of $\mathcal L$.}
For every integer $j \in \{ 1,\dots,d \}$, define $\alpha_j \in H^2(\mathbb T^{2d};\mathbb R)$ via
$$
\alpha_j:=[du_j\wedge dv_j].
$$
We now prove that
$
c_1(\mathcal L)_{\mathbb R}
=
\alpha_1+\cdots+\alpha_d.
$
To do so, we start by computing the Chern class using a connection. Although the sections
$s_\ell$ are only assumed continuous, the bundle $\mathcal L$ itself is a smooth
complex line bundle, so its Chern class may be computed using any smooth connection.
Let $\sigma:\mathbb R^{2d}\to\mathbb C$ be a smooth function satisfying
$$
\sigma(u+n,v+m)=e^{2\pi i\,n\cdot v}\sigma(u,v),
$$
and define
$$
\nabla\sigma
:=
d\sigma
-
2\pi i
\left(\sum_{j=1}^d u_j\,dv_j\right)\sigma.
$$
We check that $\nabla$ descends to a connection on $\mathcal L$. Indeed,
$$
\begin{aligned}
(\nabla\sigma)(u+n,v+m)
&=
d\!\left(e^{2\pi i\,n\cdot v}\sigma(u,v)\right)
-
2\pi i\sum_{j=1}^d (u_j+n_j)\,dv_j\,
e^{2\pi i\,n\cdot v}\sigma(u,v) \\
&=
e^{2\pi i\,n\cdot v}
\left(
d\sigma
+
2\pi i\sum_{j=1}^d n_j\,dv_j\,\sigma
-
2\pi i\sum_{j=1}^d u_j\,dv_j\,\sigma
-
2\pi i\sum_{j=1}^d n_j\,dv_j\,\sigma
\right) \\
&=
e^{2\pi i\,n\cdot v}
\left(
d\sigma
-
2\pi i\sum_{j=1}^d u_j\,dv_j\,\sigma
\right) \\
&=
e^{2\pi i\,n\cdot v}(\nabla\sigma)(u,v).
\end{aligned}
$$
Thus, $\nabla$ defines a
connection on $\mathcal L$.
The connection one-form in the above trivialization over $\mathbb R^{2d}$ is
$$
A
=
-2\pi i\sum_{j=1}^d u_j\,dv_j.
$$
Since $\mathcal L$ has rank one, the curvature is
$$
F_\nabla=dA
=
-2\pi i\sum_{j=1}^d du_j\wedge dv_j.
$$
For a complex line bundle, the real first Chern class is represented by
$
\frac{i}{2\pi}F_\nabla,
$
which implies that
$$
c_1(\mathcal L)_{\mathbb R}
=
\left[\frac{i}{2\pi}F_\nabla\right]
=
\sum_{j=1}^d [du_j\wedge dv_j]
=
\alpha_1+\cdots+\alpha_d.
$$

\smallskip
\noindent\textbf{Step 3: the top Chern class of $E_N$.}
Since
$
E_N=\mathcal L^{\oplus N},
$
the Whitney product formula gives
$
c(E_N)=c(\mathcal L)^N.
$
Because $\mathcal L$ is a line bundle,
$
c(\mathcal L)=1+c_1(\mathcal L).
$
Hence,
$
c_N(E_N)=c_1(\mathcal L)^N.
$
Passing to real cohomology and using the computation from above, we get
$$
c_N(E_N)_{\mathbb R}
=
(\alpha_1+\cdots+\alpha_d)^N.
$$
We now show that this class is nonzero. To do so, consider the $2N$-dimensional subtorus $Y$ of $\T^{2d}$ defined by
$$
Y
:=
\{u_{N+1}=\cdots=u_d=v_{N+1}=\cdots=v_d=0\}
\subset \mathbb T^{2d}.
$$
We have the properties
$
Y\cong(\mathbb T^2)^N,
$
$
\alpha_j|_Y=0
$
for $j > N$, and hence
\begin{equation}\label{eq:777}
    \left.c_N(E_N)_{\mathbb R}\right|_Y
=
(\alpha_1+\cdots+\alpha_N)^N.
\end{equation}
The cohomology ring of the torus is the exterior algebra on the degree-one classes
$[du_1],\dots,[du_d],[dv_1],\dots,[dv_d]$. In particular,
$$
\alpha_j^2
=
[du_j\wedge dv_j]\smile[du_j\wedge dv_j]
=
0.
$$
Since the classes $\alpha_j$ have degree $2$, they commute with one another. Hence
the multinomial expansion of \eqref{eq:777} gives
\begin{equation}\label{eq:888}
    (\alpha_1+\cdots+\alpha_N)^N
=
N!\,\alpha_1\smile\cdots\smile\alpha_N.
\end{equation}
Equivalently,
$$
(\alpha_1+\cdots+\alpha_N)^N
=
N!\,[du_1\wedge dv_1\wedge\cdots\wedge du_N\wedge dv_N].
$$
Giving $Y$ its standard orientation, we obtain the identity
$$
\int_Y
du_1\wedge dv_1\wedge\cdots\wedge du_N\wedge dv_N
=
1.
$$
Thus, by \eqref{eq:777} and \eqref{eq:888} we obtain
$$
\int_Y
\left.c_N(E_N)_{\mathbb R}\right|_Y
=
N!\neq 0.
$$
The latter identity implies that
$
c_N(E_N)_{\mathbb R}\neq 0.
$
In particular, we have shown that
\begin{equation}\label{eq:999}
    c_N(E_N)\neq 0
\end{equation}
in $H^{2N}(\mathbb T^{2d};\mathbb Z)$.

\smallskip
\noindent\textbf{Step 4: contradiction from a nowhere-zero section.}
By Step 1, the functions $s_1,\dots,s_N$ define a continuous section
$
S=(s_1,\dots,s_N)
$
of
$
E_N=\mathcal L^{\oplus N}.
$
By the contradiction assumption, $S$ is nowhere zero.
Since $\mathbb T^{2d}$ is a compact Hausdorff space, every complex vector bundle over
$\mathbb T^{2d}$ admits a continuous Hermitian metric. The nowhere-zero section $S$
spans a trivial complex line subbundle
$
\mathbb C S\subset E_N.
$
Taking its orthogonal complement with respect to a continuous Hermitian metric gives
a rank-$(N-1)$ complex vector bundle $E'$ such that
$
E_N\cong \underline{\mathbb C}\oplus E'.
$
Therefore, by the Whitney product formula,
$$
c(E_N)
=
c(\underline{\mathbb C})c(E')
=
c(E').
$$
Since $E'$ has complex rank $N-1$, its degree-$2N$ Chern class vanishes,
$
c_N(E')=0.
$
Hence, we have
$
c_N(E_N)=0,
$
contradicting \eqref{eq:999}.
\end{proof}

\end{document}
